% ---------------------------------------------------------------
% Preambolo per appunti universitari — skill "appunti"
% Sostituiscilo pure col tuo: servono solo gli ambienti
% defbox{titolo} ed example{titolo}.
% ---------------------------------------------------------------
\documentclass[11pt,a4paper]{article}

\usepackage[utf8]{inputenc}
\usepackage[T1]{fontenc}
\usepackage[english]{babel}
\usepackage[margin=2.5cm]{geometry}

\usepackage{amsmath,amssymb,amsthm,mathtools}
\usepackage{bm}
\usepackage{booktabs}
\usepackage{array}
\usepackage{enumitem}
\usepackage{graphicx}
\usepackage{xcolor}
\usepackage{tikz}
\usetikzlibrary{arrows.meta,positioning,shapes.geometric,fit,calc}
\usepackage[most]{tcolorbox}
\usepackage{hyperref}
\hypersetup{colorlinks=true,linkcolor=blue!50!black,citecolor=blue!50!black,urlcolor=blue!50!black}

% --- Definizioni ------------------------------------------------
\newtcolorbox{defbox}[1]{%
  enhanced, breakable,
  colback=blue!4!white, colframe=blue!55!black,
  fonttitle=\bfseries, title={#1},
  boxrule=0.6pt, arc=2pt,
  left=6pt, right=6pt, top=5pt, bottom=5pt,
  before skip=0.30cm, after skip=0.30cm
}

% --- Esempi -----------------------------------------------------
\newtcolorbox{example}[1]{%
  enhanced, breakable,
  colback=green!3!white, colframe=green!40!black,
  fonttitle=\bfseries, title={#1},
  boxrule=0.6pt, arc=2pt,
  left=6pt, right=6pt, top=5pt, bottom=5pt,
  before skip=0.30cm, after skip=0.30cm
}

% --- Riquadri opzionali ----------------------------------------
\newtcolorbox{remarkbox}[1]{%
  enhanced, breakable,
  colback=orange!4!white, colframe=orange!60!black,
  fonttitle=\bfseries, title={#1},
  boxrule=0.6pt, arc=2pt, before skip=0.30cm, after skip=0.30cm
}

% --- Scorciatoie matematiche ------------------------------------
\newcommand{\R}{\mathbb{R}}
\newcommand{\E}{\mathbb{E}}
\newcommand{\Prob}{\mathbb{P}}
\newcommand{\Var}{\operatorname{Var}}
\newcommand{\Cov}{\operatorname{Cov}}
\newcommand{\Bias}{\operatorname{Bias}}
\newcommand{\argmin}{\operatorname*{arg\,min}}
\newcommand{\argmax}{\operatorname*{arg\,max}}
\newcommand{\indep}{\perp\!\!\!\perp}
\newcommand{\norm}[1]{\left\lVert #1 \right\rVert}
\newcommand{\abs}[1]{\left\lvert #1 \right\rvert}
\newcommand{\T}{^{\top}}

\setlength{\parindent}{0pt}
\setlength{\parskip}{0.35em}
